\documentclass[aps,pre,onecolumn,10pt,superscriptaddress]{revtex4-2}
\usepackage{amsmath,booktabs}
\usepackage[hidelinks]{hyperref}
\renewcommand{\thesection}{S\arabic{section}}
\renewcommand{\thetable}{S\arabic{table}}
\renewcommand{\theequation}{S\arabic{equation}}
\begin{document}
\title{Supplemental Material: Degree Information and Norm Turnover in Random-Walk Network Embeddings}
\author{Govind Gandhi}
\affiliation{School of Data Science, University of Virginia, Charlottesville, Virginia 22904, USA}
\author{Sadamori Kojaku}
\affiliation{School of Systems Science and Industrial Engineering, Binghamton University, New York 13902, USA}
\author{Yong-Yeol Ahn}
\affiliation{School of Data Science, University of Virginia, Charlottesville, Virginia 22904, USA}
\date{\today}
\maketitle

\section{Post hoc predictor comparisons}
\label{sec:predictor-comparison}

We compare the frozen exposure--KL predictor with degree-and-threshold and single-feature baselines using the existing calibration and target records.
Every regressor uses the original ExtraTrees settings: 400 trees, minimum leaf size 40, maximum features two, and random seed 20260813.
Only configuration-model calibration outcomes enter fitting; target outcomes are used solely for evaluation.
The inputs are log positive-pair count and context KL, log degree and log threshold, or one of log pair count, context KL, and log degree alone.
These additional comparisons are post hoc, and no hyperparameters are selected from their target scores.
The original frozen predictions, protocol, and pass/fail outcomes remain unchanged.
The corrected retained-count and raw-count negative-sampling results are retrospective revalidations of previously examined targets.

Within-curve Spearman correlations compare predicted and observed individual node norms in each of nine family--threshold groups, pooling five replicates per group.
We report the median and range across these groups; the range is not a confidence interval.
The pooled correlation additionally reflects variation between families and thresholds.
Separate correlations for all 45 family--threshold--replicate groups are retained in the accompanying data.
Thresholds share graphs, and NetScience's replicates share one observed network, so these groups are not independent samples.

Table~\ref{tab:prospective-baselines} reports raw center norms, the primary response.
The degree-and-threshold baseline passes all four declared criteria in every dataset and achieves more peaks within one bin than the exposure--KL predictor in the corrected retained-noise comparison (9/9 versus 8/9).
The exposure--KL predictor instead gives higher within-curve correlations; its raw-norm correlation exceeds the degree-and-threshold baseline in all 45 separate replicate groups in each dataset.
Degree alone fails the pooled-correlation criterion, whereas pairs alone pass all four criteria in the corrected raw-noise arm.
Thus neither passing the criteria nor peak localization alone distinguishes the proposed predictor.

\begin{table}[htbp]
\centering
\small
\setlength{\tabcolsep}{6pt}
\caption{Post hoc comparisons for raw center norms (primary response). The last column applies the original four criteria, including the predictor-independent observed-peak property. Within-curve Spearman is the median [range] across nine family--threshold groups, each pooling five replicates. Corrected results are retrospective revalidations.}
\label{tab:prospective-baselines}
\begin{tabular}{lccccc}
\toprule
Inputs & \shortstack{Median peak\\error (bins)} & \shortstack{Peaks within\\one bin} & Pooled $\rho$ & Within-curve $\rho$ & All four pass \\
\midrule
\multicolumn{6}{l}{\textit{Original}} \\
Pairs + KL & 0 & 8/9 & 0.881 & $0.842 [0.506, 0.935]$ & Yes \\
Degree + threshold & 1 & 7/9 & 0.652 & $0.293 [0.163, 0.610]$ & Yes \\
Pairs only & 1 & 5/9 & 0.591 & $0.523 [0.045, 0.631]$ & No \\
KL only & 2 & 2/9 & 0.670 & $0.630 [-0.140, 0.825]$ & No \\
Degree only & 1 & 5/9 & 0.178 & $0.279 [0.123, 0.398]$ & No \\
\midrule
\multicolumn{6}{l}{\textit{Corrected, retained noise}} \\
Pairs + KL & 0 & 8/9 & 0.875 & $0.835 [0.548, 0.934]$ & Yes \\
Degree + threshold & 0 & 9/9 & 0.649 & $0.285 [0.164, 0.602]$ & Yes \\
Pairs only & 1 & 5/9 & 0.581 & $0.510 [-0.146, 0.627]$ & No \\
KL only & 2 & 2/9 & 0.666 & $0.634 [-0.089, 0.829]$ & No \\
Degree only & 1 & 6/9 & 0.169 & $0.257 [0.117, 0.382]$ & No \\
\midrule
\multicolumn{6}{l}{\textit{Corrected, raw noise}} \\
Pairs + KL & 0 & 8/9 & 0.873 & $0.817 [0.543, 0.934]$ & Yes \\
Degree + threshold & 0 & 8/9 & 0.643 & $0.284 [0.151, 0.595]$ & Yes \\
Pairs only & 1 & 6/9 & 0.593 & $0.500 [-0.015, 0.617]$ & Yes \\
KL only & 2 & 2/9 & 0.658 & $0.624 [-0.087, 0.794]$ & No \\
Degree only & 0 & 6/9 & 0.128 & $0.193 [0.041, 0.285]$ & No \\
\bottomrule
\end{tabular}
\end{table}


Table~\ref{tab:balanced-baselines} gives the balanced-norm secondary response.
The within-curve advantage persists, although the exposure--KL predictor loses to the degree-and-threshold baseline in one of the 45 separate replicate comparisons for the original balanced response; it wins all 45 in each corrected arm.
These descriptive comparisons support improved predictive ranking relative to the specified baseline, not a causal exposure--specificity mechanism or information statistically independent of degree.

\begin{table}[htbp]
\centering
\small
\setlength{\tabcolsep}{6pt}
\caption{Post hoc comparisons for balanced norms (secondary response, without a separate primary pass/fail test). Within-curve Spearman is the median [range] across nine family--threshold groups, each pooling five replicates. Corrected results are retrospective revalidations.}
\label{tab:balanced-baselines}
\begin{tabular}{lccc}
\toprule
Inputs & \shortstack{Median peak\\error (bins)} & \shortstack{Peaks within\\one bin} & Within-curve $\rho$ \\
\midrule
\multicolumn{4}{l}{\textit{Original}} \\
Pairs + KL & 0 & 8/9 & $0.840 [0.504, 0.935]$ \\
Degree + threshold & 1 & 7/9 & $0.308 [0.159, 0.614]$ \\
Pairs only & 1 & 5/9 & $0.528 [0.047, 0.623]$ \\
KL only & 2 & 2/9 & $0.617 [-0.179, 0.823]$ \\
Degree only & 1 & 5/9 & $0.276 [0.123, 0.375]$ \\
\midrule
\multicolumn{4}{l}{\textit{Corrected, retained noise}} \\
Pairs + KL & 0 & 8/9 & $0.831 [0.548, 0.934]$ \\
Degree + threshold & 1 & 9/9 & $0.306 [0.181, 0.604]$ \\
Pairs only & 1 & 5/9 & $0.529 [-0.146, 0.617]$ \\
KL only & 2 & 2/9 & $0.621 [-0.144, 0.824]$ \\
Degree only & 1 & 6/9 & $0.268 [0.103, 0.375]$ \\
\midrule
\multicolumn{4}{l}{\textit{Corrected, raw noise}} \\
Pairs + KL & 0 & 8/9 & $0.801 [0.540, 0.935]$ \\
Degree + threshold & 0 & 8/9 & $0.288 [0.144, 0.598]$ \\
Pairs only & 1 & 6/9 & $0.524 [-0.011, 0.611]$ \\
KL only & 2 & 2/9 & $0.611 [-0.140, 0.791]$ \\
Degree only & 0 & 6/9 & $0.225 [0.012, 0.347]$ \\
\bottomrule
\end{tabular}
\end{table}


\section{Interpretation of the declared criteria}
\label{sec:criteria}

Three frozen criteria evaluate predictions: median absolute peak-bin error at most one, at least two thirds of peaks within one bin, and pooled node-norm Spearman correlation at least 0.5.
The fourth asks whether observed peak degree is nondecreasing with weaker subsampling in at least two of three graph families.
It depends only on observed curves, is identical for every predictor, and permits ties; the original Barab\'asi--Albert peaks at thresholds \(10^{-3}\) and \(3\times10^{-3}\) both have degree approximately 6.77.
Keeping this response property distinct from predictive skill changes the interpretation, not the frozen evaluation.

For the next genuinely new test on natural text or mobility, comparators and evaluation rules should be specified before examining outcomes.
They should include degree or frequency baselines, threshold when varied, and pair/KL ablations, together with within-target ranking and peak-error criteria and uncertainty across independent graphs or datasets.
No new pass threshold should be chosen retrospectively from the present comparison.

\section{Native factor balance}
\label{sec:factor-balance}

\paragraph{Gradient-flow identity.}
Let \(U\) and \(V\) be real matrices with the same column dimension, \(M=UV^\top\), and \(G=\nabla_M\mathcal{L}(M)\) for a differentiable score-only loss.
Under simultaneous Euclidean gradient flow with the same rate for both factors, \(\dot U=-GV\) and \(\dot V=-G^\top U\).
Consequently,
\begin{equation}
 \frac{d}{dt}(U^\top U)=\frac{d}{dt}(V^\top V)
 =-V^\top G^\top U-U^\top GV,
 \qquad D=U^\top U-V^\top V=\text{constant}.
\end{equation}
This is the established balance conservation property for linear factors~\citep{du2018algorithmic}, not a new theorem.
When \(D=0\),
\begin{equation}
 MM^\top=U(V^\top V)U^\top=(UU^\top)^2.
\end{equation}
Uniqueness of the positive-semidefinite square root gives \(UU^\top=(MM^\top)^{1/2}\), including rank-deficient factors.
Its diagonal gives equality of raw and canonical balanced squared row norms.
Conservation does not drive a nonzero \(D\) to zero: with \(V=0\) and small nonzero \(U\), absolute imbalance is small but \(\|D\|_F/\|U^\top U\|_F=1\).

For one simultaneous gradient-descent step of size \(\eta\), with both gradients evaluated at the pre-update factors,
\begin{equation}
 D_{\rm next}-D=\eta^2\left(V^\top G^\top GV-U^\top GG^\top U\right).
\end{equation}
The first-order terms cancel; the remaining terms need not vanish.
This formula also applies to a score-only sampled loss with simultaneous updates, but not verbatim to Gensim's sequential updates.
Unequal factor learning rates, adaptive updates, or factor-specific regularization are outside the stated flow identity.
Neither balance nor norm agreement establishes score convergence or identifies a norm--degree profile independently of training.

\paragraph{Native training diagnostic.}
We replay all twenty corrected Gensim controls: five graph seeds at each of four filtering thresholds, at the original five-epoch endpoints.
Graphs, walks, vocabulary and initialization signatures, training-token counts, and both final factor hashes match the recorded runs exactly.
We measure the native factors before any balancing transformation.
Table~\ref{tab:factor-balance} compares the requested relative Gram imbalance with actual norm discrepancies, where \(r_i=\|u_i\|\) and \(r_{{\rm bal},i}\) is the canonical balanced norm.
The individual-node error is \(|r_i-r_{{\rm bal},i}|/r_i\).

\begin{table}[htbp]
\caption{Native Gensim factor balance at the original five-epoch endpoints. Ranges span five graph seeds at each threshold; the last column is the largest individual-node discrepancy across those runs.}
\label{tab:factor-balance}
\centering
\begin{tabular}{lrrr}
\toprule
Threshold & $\|D\|_F/\|U^\top U\|_F$ & $\|r-r_{\rm bal}\|_2/\|r\|_2$ (\%) & Max. node error (\%) \\
\midrule
No filtering & 0.1347--0.1630 & 1.054--1.150 & 12.45 \\
0.0003 & 0.0786--0.0803 & 0.548--0.567 & 1.88 \\
0.001 & 0.0944--0.1066 & 0.648--0.726 & 3.48 \\
0.003 & 0.1264--0.1491 & 0.874--0.968 & 6.68 \\
\bottomrule
\end{tabular}
\end{table}


Final relative Gram imbalance ranges from \(0.0786\) to \(0.163\); it is not zero.
Absolute imbalance grows from \(0.407\)--\(0.411\) initially to approximately \(225\)--\(910\), so these finite-step runs do not approximately conserve the initial absolute imbalance.
Nevertheless, relative norm-vector discrepancies are \(0.55\)--\(1.15\%\), with a maximum individual discrepancy of \(12.45\%\).
Thus close aggregate norm agreement is an observed property of these runs, not a quantitative consequence of conserved initialization imbalance.
The diagnostic does not cover the custom mini-batch pipeline or later learning schedules; already balanced factors would not test native balance.

\paragraph{Reproduction.}
Run \texttt{python3}~\path{scripts/factor_balance.py} using the pinned, patched runtime documented in \path{analysis/negative_sampling_control_protocol.md}.
The script checks source hashes, dependency versions, and the source completion manifest before replaying the existing conditions; it generates Table~\ref{tab:factor-balance} from the same records as \path{results/factor_balance/summary.json}.
The latter retains initial and final diagnostics and raw and balanced norms for every node in every run, with final factor hashes.
No original result or training configuration is changed.


\section{Reproducibility}

The repository script \path{scripts/prospective_baselines.py} verifies the original freeze hashes, corrected completion manifest, dependency versions, and target-record alignment before fitting baselines.
It reproduces the frozen predictor exactly and generates both tables from the same summaries.
All per-curve and per-replicate records, including failures, are in \path{results/prospective_baselines/summary.json}; setup and commands are in \path{analysis/prospective_baselines.md}.
No new embeddings are trained for these comparisons.
The separate factor-balance audit deterministically replays existing training runs, as described in Sec.~\ref{sec:factor-balance}.

\bibliography{references}
\end{document}
